\section*{結言}
今回の8件では、資本市場への展開と実務向けモデルの競争が並行して進んだ。AnthropicのS-1提出は事業規模や投資負担を公開資料で検討する段階への移行を示し、GPT-6.1 TerraやCursorのDeepSeek対応は、モデルの利用料金と作業効率を巡る比較を加速させている。ただし、価格性能の優位性は発表者の評価条件を踏まえて読む必要がある。

Cohereの報告は、モデルが説明した意図と実際のツール操作が一致するとは限らないという運用上の論点を提示した。CodeRabbitやShannonの話題も、開発現場にAIを組み込む際の検証と安全性の重要性につながる。一方、Shannonの実績やGrok Botの更新内容には、公式な裏付けが不足している部分が残る。

今後は、S-1の財務注記、Terraの独立評価、Cohereの技術報告と再現性、各ツールの正式な仕様が注目点となる。なお、CursorとUnslothの投稿は観測窓の終端後であり、本号では時刻の例外を明記して紹介した。
